\title{QLoRA: Efficient Finetuning of Quantized LLMs}

\begin{abstract}
We introduce \textsc{QLoRA}, an effective finetuning technique designed to minimize memory requirements, enabling the finetuning of a 65B parameter model on a single 48GB GPU while achieving performance comparable to traditional 16-bit finetuning. In \textsc{QLoRA}, gradients are propagated through a fixed, 4-bit quantized pretrained language model into Low Rank Adapters~(LoRA). Our premier model suite, termed \textbf{Guanaco}, surpasses all previous open-source models on the Vicuna benchmark, attaining 99.3\% of ChatGPT’s performance after just 24 hours of finetuning on a single GPU. Key innovations in \textsc{QLoRA} enabling memory efficiency without compromising on accuracy include: (a) 4-bit NormalFloat (NF4), a novel datatype theoretically optimized for normally distributed weights, (b) Double Quantization, which compresses memory use further by quantizing the quantization factors themselves, and (c) Paged Optimizers to prevent memory usage peaks. Utilizing \textsc{QLoRA}, we have finetuned over 1,000 models, offering a comprehensive assessment of instruction-following and chatbot effectiveness across 8 instructional datasets, various architectures (LLaMA, T5), and large model scales otherwise impractical to finetune (such as 33B and 65B parameters). Our experiments reveal that finetuning with QLoRA on smaller, high-quality datasets achieves cutting-edge results—even with models smaller than prior leading approaches. Additionally, we present an in-depth evaluation of chatbot capabilities through both human and GPT-4 assessments, finding that GPT-4 based evaluation provides a low-cost and sufficiently reliable proxy for human judgment. Notably, our analysis also highlights the unreliability of current chatbot benchmarks in accurately gauging model quality. We further perform a targeted “lemon-picked” analysis highlighting areas where \textbf{Guanaco} underperforms relative to ChatGPT. All models and source code, including CUDA kernels for 4-bit training, are made publicly available.\footnote{\url{https://github.com/artidoro/qlora} and \url{https://github.com/TimDettmers/bitsandbytes}}
\end{abstract}

\section{Introduction}

Customizing large language models (LLMs) through finetuning is a proven approach for enhancing their capabilities \citep{min2021metaicl, wei2021finetuned, ouyang2022training, wang2022super, wang2022self, liu2022few}, as well as for instilling or suppressing particular model behaviors \citep{ouyang2022training,askell2021general,bai2022training}. Nevertheless, finetuning state-of-the-art, large-scale models remains highly resource-intensive; for example, standard 16-bit finetuning of the LLaMA model with 65B parameters~\citep{touvron2023llama} demands upwards of 780 GB of GPU memory. Although emerging quantization techniques offer a means of lowering the inference memory demands for LLMs \citep{dettmers2022llm, dettmers2022case, frantar2022gptq, xiao2022smoothquant}, such methods typically prove inadequate for training scenarios, where their benefits diminish or fail entirely \citep{wortsman2023stable}.

In this work, we provide the first evidence that finetuning a 4-bit quantized model is achievable without sacrificing model performance. Our approach, called \textsc{QLoRA}, introduces an advanced quantization mechanism that converts pretrained models to 4-bit precision while incorporating a compact set of trainable Low-rank Adapter weights \citep{hu2021lora}. 

These adapter weights are optimized by propagating gradients through the quantized parameters.

Employing \textsc{QLoRA} enables a reduction of average GPU memory usage for finetuning a 65B model from over 780GB to under 48GB, while maintaining equivalent runtime and accuracy relative to conventional 16-bit finetuning. This innovation dramatically broadens the range of users who can undertake large-scale LLM finetuning, allowing even the largest open-source models to be finetuned on a single GPU. Leveraging \textsc{QLoRA}, we introduce the \textbf{Guanaco} suite of models, where the second most performant member achieves 97.8\% of ChatGPT's score on the Vicuna~\citep{vicuna2023} benchmark and can be trained in less than 12 hours on a standard consumer GPU. With a single professional GPU over 24 hours, our largest Guanaco model attains 99.3\% of ChatGPT's performance, virtually matching ChatGPT on this benchmark. Furthermore, at deployment, our smallest \textbf{Guanaco} model (with 7B parameters) operates with a mere 5 GB memory requirement and outperforms a 26 GB Alpaca model by over 20 percentage points on the Vicuna benchmark (see Table~\ref{tbl:vicuna_eval}).

\textsc{QLoRA} brings forth several advancements aimed at minimizing memory consumption while maintaining robust performance: (1) \textbf{4-bit NormalFloat}, a quantization datatype theoretically optimized for data following a normal distribution, which outperforms both 4-bit Integers and 4-bit Floats in empirical evaluations. (2) \textbf{Double Quantization}, which further compresses model size by quantizing the quantization coefficients themselves, achieving an average reduction of approximately 0.37 bits per parameter (equivalent to roughly 3 GB in a 65B parameter model). (3) \textbf{Paged Optimizers}, leveraging NVIDIA's unified memory framework to circumvent transient memory spikes associated with gradient checkpointing during processing of mini-batches with extended sequence lengths. By integrating these innovations, we refine the LoRA methodology by embedding adapters throughout every network layer, thereby nearly eliminating the accuracy compromises that have characterized earlier approaches.

The memory efficiency offered by \textsc{QLoRA} allows for comprehensive experimentation with instruction finetuning and chatbot evaluation across model scales that conventional finetuning would render impractical due to prohibitive memory demands. Utilizing this approach, we train upwards of 1,000 models spanning multiple instruction tuning datasets, diverse model architectures, and scales ranging from 80M to 65B parameters. Beyond demonstrating that \textsc{QLoRA} matches the performance of 16-bit finetuning (\S\ref{sec:comp_to_std}) and facilitating the development of a state-of-the-art chatbot, \textbf{Guanaco} (\S\ref{sec:pushing}), our study delves into observed patterns across trained models. Notably, our findings indicate that dataset quality significantly outweighs dataset size; for example, the 9k sample OASST1 dataset surpasses the performance of a substantially larger 450k sample FLAN v2 dataset (when subsampled) in the context of chatbot instruction following. Additionally, we reveal that excelling on the Massive Multitask Language Understanding (MMLU) benchmark does not guarantee strong outcomes on the Vicuna chatbot benchmark and vice versa, highlighting that task-appropriate data curation is more critical than sheer dataset size.

In addition, we conduct a comprehensive assessment of chatbot systems utilizing both human evaluators and GPT-4 as judges. Our evaluation framework adopts a tournament approach, where different models are pitted against each other, each generating a response to the same prompt, with the superior reply advancing as determined either by GPT-4 or human raters. These match outcomes are synthesized into Elo scores~\citep{elo1967proposed, elo1978rating}, which provide an ordinal ranking of chatbot capabilities. Our experiments indicate that there is broad concordance between GPT-4-based and human-driven rankings, though notable instances of divergence are also observed. We emphasize, therefore, that although model-based evaluation can offer a scalable and low-cost alternative to manual annotation, it also introduces its own sources of variability.

To supplement the quantitative benchmarking, we present a qualitative examination of \textbf{Guanaco} models. This analysis surfaces both strengths and shortcomings in model behavior not identified through numerical metrics alone.

For reproducibility and to foster future research, we make publicly available all generated model outputs with both human and GPT-4 judgments. Our software and CUDA implementations are released as open source, with methods integrated into the Hugging Face transformers library~\citep{wolf2019huggingface}, ensuring broad accessibility. We also provide a suite of adapters compatible with 7/13/33/65B parameter models, trained on eight diverse instruction-following datasets, collectively yielding 32 distinct, publicly available, fine-tuned models.

\section{Background}
\paragraph{Block-wise k-bit Quantization} Quantization refers to the conversion of data from a high-precision format to one of reduced precision, typically by transforming data types with greater bit-widths to types with fewer bits—such as converting 32-bit floating-point values into 8-bit integer values. To optimally exploit the available range in the lower-precision data type, it is customary to normalize the original data, usually a tensor, by its maximum absolute value before discretization. As an illustration, when mapping an FP32 tensor into an Int8 tensor with permissible values in $[-127, 127]$:

\begin{equation}
    \mathbf{X}^{ \text{Int8}} = \text{round}\left( \frac{127}{{\text{absmax}}(\mathbf{X}^{{\text{FP32}}})}\mathbf{X}^{{\text{FP32}}} \right) = \text{round}(c^{\text{FP32}}\cdot \mathbf{X}^{{\text{FP32}}}),
\end{equation}

where $c$ is the {\em quantization constant} or {\em quantization scale}. The corresponding dequantization operation is defined as:

\begin{equation}
    \text{dequant}(c^{\text{FP32}}, \mathbf{X}^{\text{Int8}}) = \frac{\mathbf{X}^{\text{Int8}}}{c^{\text{FP32}}} = \mathbf{X}^{\text{FP32}}
\end{equation}

A limitation of this method is its susceptibility to outliers; if the input tensor contains values with large magnitudes, quantization bins—certain bit patterns—may become underutilized or empty, as very few elements are mapped to them. To mitigate this issue, a widely adopted strategy is to divide the input tensor into several blocks and quantize each block separately, assigning an independent quantization constant $c$ to each one. More concretely, the input tensor $\mathbf{X}\in \mathbb{R}^{b\times h}$ is first flattened and then partitioned into $n$ consecutive segments, each of size $B$, leading to ${n = ({b\times h} )/{B} }$ blocks in total. Quantization is performed independently within each block according to Equation~1, resulting in a quantized representation along with $n$ associated quantization constants $c_i$.

\paragraph{Low-rank Adapters} The Low-rank Adapter (LoRA) approach for finetuning \citep{hu2021lora} addresses memory constraints by introducing a compact set of tunable adapter parameters, while keeping the majority of model weights fixed and immutable. During training via stochastic gradient descent, gradients are computed through the static pretrained parameters but parameter updates are restricted to the adapter components, which are trained to reduce the loss. LoRA enhances a linear layer by adding a low-rank, factorized parameterization. Given a linear transformation ${\mathbf{X} \mathbf{W} = \mathbf{Y}}$, where $\mathbf{X}\in  \mathbb{R}^{b\times h}$ and $\mathbf{W}\in  \mathbb{R}^{h\times o}$, LoRA modifies this projection as follows:

\begin{equation}
    \mathbf{Y} = \mathbf{X}\mathbf{W} + s\mathbf{X}\mathbf{L}_1\mathbf{L}_2,
\end{equation}

in which $\mathbf{L}_1\in  \mathbb{R}^{h\times r}$, $\mathbf{L}_2\in  \mathbb{R}^{r\times o}$, and $s$ is a scaling factor.

\paragraph{Memory Requirement of Parameter-Efficient Finetuning} An important aspect to consider is the memory consumption of LoRA during training, particularly regarding both the quantity and dimensions of adapters deployed. Given that LoRA imposes a very small memory overhead, it becomes feasible to incorporate additional adapters to enhance performance, with negligible impact on total memory consumption. Although LoRA was introduced as a Parameter Efficient Finetuning (PEFT) strategy, the primary source of memory usage in LLM finetuning actually stems from the storage of activation gradients rather than from LoRA's learnable parameters. For example, finetuning a 7B LLaMA model on FLAN v2 using a batch size of 1, where the LoRA parameters constitute the typical 0.2\% of the original model's weights\citep{hu2021lora,liu2022few}, results in LoRA input gradients occupying 567 MB of memory, whereas the actual LoRA parameters themselves require only 26 MB. Applying gradient checkpointing~\citep{chen2016training} brings the memory used by input gradients down to about 18 MB per sequence, still making gradients more demanding than the LoRA parameters as far as memory is concerned. By contrast, the memory usage for a 4-bit quantized base model is 5,048 MB. These observations underscore the significance of gradient checkpointing, and suggest that further reducing the size of LoRA parameters will only slightly affect memory requirements. Consequently, utilizing more adapters is possible without a considerable rise in total training memory usage (refer to Appendix~\ref{app:memory} for an in-depth analysis). This capability, as explored in later sections, is key to attaining full 16-bit precision performance.

\section{\textsc{QLoRA} Finetuning}

\textsc{QLoRA} enables accurate 4-bit finetuning through two novel approaches introduced in this work: 4-bit NormalFloat (NF4) quantization and Double Quantization. To address the issue of memory spikes during gradient checkpointing—which have historically limited the ability to finetune large models on a single device—we also present Paged Optimizers as a remedy to out-of-memory failures.

In \textsc{QLoRA}, a single low-precision format is used for storage (typically 4-bit), paired with a computational data type, most commonly BFloat16. Operationally, this means that before using any \textsc{QLoRA} weight tensor, it is first dequantized to BFloat16 and then matrix multiplications are performed in 16-bit precision.

We proceed by elaborating on the main components of \textsc{QLoRA}, followed by a formal characterization of the approach.

\paragraph{4-bit NormalFloat Quantization} The NormalFloat (NF) representation extends the principle of Quantile Quantization \citep{dettmers2022optimizers}, a theoretically optimal scheme in which each bin captures an equal proportion of elements from the input tensor. This is accomplished by approximating the input tensor's quantiles using the empirical cumulative distribution function.

However, quantile quantization suffers from a significant drawback: accurate estimation of quantiles is computationally intensive. To overcome this, efficient algorithms like SRAM quantiles~\citep{dettmers2022optimizers} are employed. Nevertheless, the approximation inherent in such algorithms introduces substantial quantization errors for outlier values, which can be among the most critical in the tensor.

These computational and accuracy issues become negligible if input tensors adhere to a distribution determined up to a quantization constant. Under such circumstances, tensors share the same quantile boundaries, making precise quantile computations feasible.

Pretrained neural network weights typically follow a zero-mean normal distribution characterized by standard deviation $\sigma$ (refer to Appendix~\ref{app:normality}). By adjusting $\sigma$, it is possible to transform all such weights so that their distribution precisely spans the representational limits of the chosen data type. Here, we assign the interval $[-1, 1]$ as the operational range for our data type. Consequently, both the weight values and the data type quantiles must be normalized to this interval.

For zero-mean normal distributions with arbitrary $\sigma$, confined within $[-1, 1]$, the most information-efficient data type can be established in the following way: (1) compute the $2^k+1$ quantiles of a standard $N(0, 1)$ distribution, yielding a $k$-bit quantile-based quantization scheme for normal distributions, (2) map the resulting quantile values into the $[-1, 1]$ interval, and (3) quantize any input weight tensor by bringing its values into $[-1, 1]$ through rescaling based on the absolute maximum.

Once the scale of the weight tensor aligns with the data type's range, standard quantization can proceed. In practice, step (3) effectively modifies the standard deviation of the tensor so that it corresponds to the standard deviation inherent to the $k$-bit data type. Formally, the $2^k$ data type values $q_i$ can be expressed as:

\begin{equation}
q_i  = \frac{1}{2}\left( Q_X\left(\frac{i}{2^k+1}\right) + Q_X\left(\frac{i+1}{2^k+1}\right)\right),
\end{equation}

where $Q_X(\cdot)$ denotes the quantile function for the standard normal distribution $N(0,1)$. A limitation of the symmetric k-bit quantization scheme is its inability to exactly represent zero, a critical requirement for accurately quantizing padding values and other inherently zero-valued entries without incurring error. To guarantee the presence of a discrete zeropoint at $0$ and to utilize the full spectrum of $2^k$ representable values in a k-bit format, we construct an asymmetric data type. This involves determining quantiles $q_i$ over two segments: $2^{k-1}$ for the negative domain and $2^{k-1}+1$ for the positive domain. After computing these quantile sets, we merge them and eliminate the redundant zero that appears in both portions. We refer to this asymmetric type, which guarantees equal expected occupancy of each quantization bin, as the \emph{k-bit NormalFloat} (NFk) since, for zero-mean normal data, it achieves information-theoretic optimality. Further details on the precise quantized values can be found in Appendix~\ref{app:nf4}.

\paragraph{Double Quantization{}} 
We propose \emph{Double Quantization} (DQ) as a strategy that compresses the quantization constants themselves to yield additional memory efficiency. Achieving accurate 4-bit quantization typically necessitates small blocks~\citep{dettmers2022case}, yet this introduces substantial overhead from storing quantization constants. As an illustration, when 32-bit quantization constants are used for a block size of 64 in $\mathbf{W}$, the constants themselves account for $32/64 = 0.5$ bits per parameter, significantly increasing memory use. Double Quantization effectively lessens this overhead.

In more detail, Double Quantization considers the initial set of quantization constants, $c_2^{\text{FP32}}$, as input for a second quantization operation. The result of this step is a set of quantized quantization constants, $c_2^{\text{FP8}}$, along with a new set of quantization constants at this higher quantization level, $c_1^{\text{FP32}}$. For this secondary quantization stage, we employ 8-bit floating point numbers and a block size of 256, as previous findings~\citep{dettmers2022case} show that 8-bit quantization induces no perceptible loss in accuracy. Given that the original $c_2^{\text{FP32}}$ values are strictly non-negative, we first mean-center them (by subtracting their mean) to make them suitable for symmetric quantization. This approach, assuming a block size of 64, lowers the per-parameter storage from $32/64=0.5$ bits to $8/64 + 32/(64\cdot256) = 0.127$ bits, achieving a per-parameter saving of 0.373 bits.

\paragraph{Paged Optimizers} leverage NVIDIA's unified memory~\footnote{\tiny \url{https://docs.nvidia.com/cuda/cuda-c-programming-guide}}, which transparently manages page-by-page data movement between CPU and GPU, thereby preventing errors when GPU memory becomes insufficient. This mechanism resembles standard operating system memory paging, which swaps data between system RAM and disk storage as required. Within our approach, optimizer state memory is allocated as paged, allowing optimizer state data to be offloaded automatically to the host CPU memory upon GPU memory exhaustion, and later restored to GPU memory during the optimizer's update operations as necessary.

\paragraph{\textsc{QLoRA}.} Building upon the previously described methods, we formalize \textsc{QLoRA} for a quantized base model's single linear layer augmented with one LoRA adapter as follows:

\begin{equation}
    \mathbf{Y}^{\text{BF16}} =  \mathbf{X}^{\text{BF16}} \text{doubleDequant}(c_1^{\text{FP32}}, c_2^{\text{k-bit}}, \mathbf{W}^{\text{NF4}}) + \mathbf{X}^{\text{BF16}}\mathbf{L}^{\text{BF16}}_1\mathbf{L}^{\text{BF16}}_2,
\end{equation}

where the \textit{doubleDequant} operation is specified by:

\begin{equation}
    \text{doubleDequant}(c_1^{\text{FP32}}, c_2^{\text{k-bit}}, \mathbf{W}^{\text{k-bit}}) = \text{dequant}(\text{dequant}(c_1^{\text{FP32}}, c_2^{\text{k-bit}}), \mathbf{W}^{\text{4bit}}) = \mathbf{W}^{\text{BF16}},
\end{equation}

For these computations, $\mathbf{W}$ is encoded in NF4 format, and $c_2$ utilizes FP8 precision.

To optimize for quantization accuracy, a blocksize of 64 is selected for $\mathbf{W}$, while for $c_2$, a blocksize of 256 is adopted to optimize memory usage.

During parameter optimization, only the gradients with respect to the adapter matrices, $\frac{\partial E}{\partial \mathbf{L}_i}$, are necessary; gradients for the 4-bit weight matrix, $\frac{\partial E}{\partial \mathbf{W}}$, are not computed. Nevertheless, determining $\frac{\partial E}{\partial \mathbf{L}_i}$ still requires the evaluation of $\frac{\partial \mathbf{X}}{\partial \mathbf{W}}$, which is derived by following equation~(5): the weights $\mathbf{W}^{\text{NF4}}$ are dequantized to $\mathbf{W}^{\text{BF16}}$ for calculation, ensuring the derivative $\frac{\partial \mathbf{X}}{\partial \mathbf{W}}$ is determined using BFloat16 arithmetic.

In summary, \textsc{QLoRA} employs a single data type for storage, typically 4-bit NormalFloat, alongside a distinct data type for computation, 16-bit BrainFloat. Prior to both the forward and backward computations, stored values are dequantized from the 4-bit representation into 16-bit BrainFloat. However, gradient computations are restricted to the LoRA weights, which remain in the 16-bit BrainFloat format.

\section{QLoRA vs. Standard Finetuning}
\label{sec:comp_to_std}

We have provided an overview of QLoRA's mechanism and outlined how it enables substantial memory savings during model finetuning. A central issue that arises is whether QLoRA can achieve comparable results to conventional full-model finetuning. Moreover, we are interested in disentangling the contributions of various QLoRA components, particularly examining how NormalFloat4 fares against the more traditional Float4. The experiments described in the sections below are designed to address these points.

\paragraph{Experimental setup.} Our investigation considers three types of architectures: encoder, encoder-decoder, and decoder-only models. QLoRA is compared both to 16-bit adapter-finetuning and to complete finetuning for models as large as 3B parameters. The evaluation encompasses benchmarks such as GLUE \citep{wang2018glue} for RoBERTa-large \citep{liu2019roberta}, Super-NaturalInstructions (TKInstruct) \citep{wang2022super} using T5 \citep{t5}, and 5-shot MMLU \citep{hendrycksmeasuring} following LLaMA finetuning on Flan v2 \citep{longpre2023flan} and Alpaca \citep{alpaca}. To further investigate NF4’s merits relative to other 4-bit quantization types, we replicate the setup from \citet{dettmers2022case}, assessing zero-shot post-quantization performance (accuracy and perplexity) for a range of models (OPT~\citep{zhang2022opt}, LLaMA~\citep{touvron2023llama}, BLOOM~\citep{scao2022bloom}, and Pythia~\citep{biderman2023pythia}) spanning model sizes from 125m to 13B parameters.

Detailed explanations are included within each result subsection for clarity, and the Appendix~\ref{app:exp-setup} contains comprehensive methodological information.

Paged optimizers are essential for \textsc{QLoRA} tuning with 33B/65B models using a single 24/48GB GPU; however, quantitative results for Paged Optimizers are omitted because paging events are infrequent, mainly triggered by mini-batches containing very long sequences. Instead, we analyze runtime performance for paged optimizers on 65B models with 48GB GPUs and observe that, for a batch size of 16, these optimizers do not reduce training speed compared to standard optimizers. Future research should systematically identify when paging-induced slowdowns might manifest.

\paragraph{Default LoRA hyperparameters do not match 16-bit performance} 

Applying LoRA only to the query and value attention projection matrices \citep{hu2021lora}—the conventional approach—does not achieve performance parity with full-model finetuning for large-scale models. Figure~\ref{fig:hyper} (LLaMA 7B on Alpaca) illustrates that the total number of LoRA adapters is the dominant hyperparameter affecting performance; applying LoRA to every linear layer within the transformer block is necessary for equivalency with full finetuning. Other hyperparameters, such as the rank $r$ of the projection, have negligible impact (see Appendix \ref{app:exp-setup} for details).

In a similar vein, we observe that the default hyperparameters used for fully finetuned baseline models are not sufficiently optimized. To establish strong baselines, we conduct a hyperparameter sweep, exploring learning rates in the range of $1 \times 10^{-6}$ to $5 \times 10^{-5}$ and batch sizes spanning from 8 to 128. The outcomes for finetuning the 7B LLaMA model on Alpaca can be seen in Figure~\ref{fig:hyper}.

\paragraph{4-bit NormalFloat surpasses 4-bit Floating Point in performance} 

Although the 4-bit NormalFloat (NF4) data format is theoretically optimal from an information perspective, it is essential to evaluate whether this theoretical benefit leads to practical improvements. Following the protocol described by \citet{dettmers2022case}, we assess quantized large language models (OPT \citep{zhang2022opt}, BLOOM \cite{scao2022bloom}, Pythia \cite{biderman2023pythia}, LLaMA) across various model sizes (from 125M to 65B parameters) and datatype configurations, focusing on both language modeling and a suite of zero-shot benchmarks. As depicted in Figure~\ref{fig:nf4} and Table~\ref{tbl:nf4_ppl}, our results indicate that NF4 offers substantial improvements over FP4 and Int4 formats, while employing double quantization further decreases memory usage without impairing model performance.

\paragraph{k-bit \textsc{QLoRA} achieves parity with 16-bit full finetuning and 16-bit LoRA} Prior work has demonstrated the feasibility of applying 4-bit quantization for \emph{inference}, albeit at the cost of diminished performance compared to 16-bit approaches \citep{dettmers2022case,frantar2022gptq}. This prompts the important question of whether the observed performance loss can be restored via 4-bit adapter finetuning. We address this using two experimental scenarios.

In the first, we perform a direct comparison with full 16-bit finetuning on RoBERTA and T5 models with parameter counts between 125M and 3B, evaluated on both the GLUE and Super-NaturalInstructions datasets. As detailed in Table~\ref{tab:nf4vsbf16}, we find that adapter-based methods using 16-bit, 8-bit, and 4-bit quantization all attain results that match those of the 16-bit fully finetuned baseline. These observations imply that any performance reduction due to low-precision quantization can be fully mitigated by subsequent adapter finetuning.

In our second experimental configuration, given that fully finetuning models with 11B parameters or more surpasses the memory capacity of a single high-memory GPU server, we further investigate whether 4-bit \textsc{QLoRA} delivers performance on par with 16-bit LoRA across model sizes from 7B to 65B parameters. Specifically, we conduct finetuning on LLaMA models ranging from 7B to 65B parameters using the Alpaca and FLAN v2 instruction-following datasets, and subsequently assess the models using the MMLU benchmark under 5-shot accuracy conditions. The outcomes, presented in Table~\ref{tbl:llama-datatype}, demonstrate that the NF4 quantization scheme with double quantization accurately restores the MMLU performance observed with 16-bit LoRA. Additionally, our analysis reveals that \textsc{QLoRA} implemented with FP4 quantization underperforms compared to the 16-bit brain float LoRA baseline by nearly 1 percentage point. These findings reinforce two key observations: (1) \textsc{QLoRA} utilizing NF4 matches the results of both 16-bit full finetuning and 16-bit LoRA finetuning, and (2) NF4 provides greater quantization precision compared to FP4.

\paragraph{Summary} Across a range of academic evaluation benchmarks, our empirical evidence consistently demonstrates that 4-bit \textsc{QLoRA} with the NF4 datatype achieves parity with 16-bit full finetuning and 16-bit LoRA approaches. Our results further establish the superiority of NF4 over FP4 and confirm that the addition of double quantization maintains, rather than diminishes, performance levels. Taken together, these findings substantiate that 4-bit \textsc{QLoRA} is a robust approach, matching the effectiveness of traditional 16-bit methods.

Building upon insights from prior work on quantization \citep{dettmers2022case}, our MMLU and Elo evaluations suggest that, for a given resource constraint in finetuning and inference, scaling up the parameter count of the base model while using lower-precision representations can be advantageous. This underscores the efficiency gains provided by \textsc{QLoRA}. As our experiments with 4-bit finetuning reveal no loss in performance relative to full-finetuning, a natural direction for future investigation is to delineate the precise boundary where the performance-precision trade-off occurs for \textsc{QLoRA} finetuning.

We undertake the exploration of instruction tuning at scales unattainable with conventional 16-bit finetuning methods given the hardware constraints typical of academic research environments.

\section{Advancing Chatbot Capabilities with QLoRA}
\label{sec:pushing}
Having demonstrated that 4-bit \textsc{QLoRA} achieves parity with traditional 16-bit performance across various scales, datasets, and tasks, we extend our investigation to comprehensive instruction finetuning using the largest open-source language models accessible to the research community. Our assessment encompasses both the challenging MMLU benchmark for Natural Language Understanding and the design of novel evaluation approaches for real-world chatbot efficacy.

\subsection{Experimental Setup} Here, we summarize our experimental design; comprehensive methodological specifics are provided in Appendix \ref{app:chatbot}.

\paragraph{Data} Noting the absence of a systematic analysis of recently introduced instruction-following corpora, we curate a collection of eight contemporary datasets. Our selection spans crowd-sourced data (OASST1~\citep{kopf2023openassistant}, HH-RLHF~\citep{bai2022training}), instruction-tuned model distillations (Alpaca~\citep{alpaca}, self-instruct~\citep{wang2022self}, unnatural-instructions~\citep{honovich2022unnatural}), large-scale corpus assemblies (FLAN v2~\citep{chung2022scaling}), as well as mixed-type datasets (Chip2~\citep{laion2023}, Longform~\citep{koksal2023longform}). Collectively, these datasets offer broad coverage in terms of linguistic diversity, dataset scale, and licensing frameworks.

\paragraph{Training Setup} To ensure results are not confounded by differing training protocols, all QLoRA finetuning is performed using cross-entropy loss under supervised learning, omitting reinforcement learning approaches—even for datasets with human evaluation scores for multiple candidate responses. In cases where datasets delineate between prompts and answers, only the responses are used for finetuning (see Appendix \ref{app:chatbot} for related ablation studies). When multiple response options exist, as in OASST1 and HH-RLHF, we systematically select the most highly rated response at each conversational juncture and train on the complete selected dialogue, including instruction segments. Across experiments, we consistently utilize NF4 \textsc{QLoRA} configurations, integrating double quantization and paged optimizers to manage memory consumption spikes associated with gradient checkpointing. Limited hyperparameter optimization is undertaken for 13B and 33B LLaMA architectures, with findings indicating that configurations established at 7B are broadly transferable—learning rate and batch size being exceptions. For 33B and 65B models, the learning rate is halved while the batch size is doubled.

\paragraph{Baselines} Model performance is benchmarked against both research (Vicuna~\citep{vicuna2023}, Open Assistant~\citep{kopf2023openassistant}) and proprietary commercial systems (GPT-4 \citep{openai2023gpt}, GPT-3.5-turbo, Bard). Specifically, Open Assistant constitutes a LLaMA 33B model trained with RLHF on OASST1, aligning with one of our primary datasets. Vicuna, by contrast, is derived via full fine-tuning of LLaMA 13B using proprietary conversational data sourced from ShareGPT, representing a distilled variant of OpenAI GPT model outputs.

\subsection{Evaluation}

Adhering to established protocols, we utilize the MMLU (Massively Multitask Language Understanding) benchmark \citep{hendrycksmeasuring} to evaluate model competence over a broad suite of language understanding challenges. This benchmark spans 57 distinct topics, ranging from elementary mathematics and US history to computer science and law. We report model performance based on 5-shot test accuracy.

Generative language abilities are further assessed using both machine-based and human-centered evaluation methods. This secondary set of assessments is grounded in human-curated queries, designed to evaluate the quality of model-generated responses. Although this type of testing provides a more authentic measure of chatbot capabilities and is increasingly prevalent, a standard protocol is yet to emerge in scholarly work. Below, we detail our methodology, which consistently employs nucleus sampling with $p=0.9$ and a temperature setting of $0.7$.

\paragraph{Benchmark Data} Two tailored datasets form the basis for our query evaluations: the Vicuna prompts \citep{vicuna2023} and the OASST1 validation set \citep{kopf2023openassistant}. The Vicuna prompts consist of 80 diverse questions utilized in their original form. In contrast, the OASST1 dataset comprises multilingual, crowd-sourced dialogues spanning multiple turns between a user and an assistant. For this benchmark, all user messages from the validation set are extracted as queries, incorporating prior conversation turns within the prompt. This process yields a total of 953 distinct user queries. These two resources are referenced as the Vicuna and OA benchmarks in our analyses.

\paragraph{Automated Evaluation} To start, in line with the protocol presented by \citet{vicuna2023}, we leverage GPT-4 for system rating on the Vicuna benchmark, comparing candidate systems to ChatGPT (GPT-3.5 Turbo). For each query, GPT-4 reviews both ChatGPT's and the system's responses, assigns each a score out of ten, and justifies its decision. The aggregate model performance is then reported as a percentage of ChatGPT's score, which may surpass 100\% should the model earn a higher score. We identify a substantial position bias whereby GPT-4 favors the first response in the prompt. To address this, we advocate averaging results across both response orders.

Subsequently, we assess performance through side-by-side output comparisons. The evaluation scheme is streamlined to a tri-class system to accommodate possible ties. GPT-4 is instructed to select the superior response, declare a tie, and provide an explanation. Comprehensive pairwise comparisons are performed between all model pairs across both the Vicuna and OA benchmarks.

\paragraph{Human Evaluation} Although contemporary studies suggest large generative models are suitable for system evaluation tasks \citep{fu2023gptscore}, there remains insufficient evidence connecting GPT-4 assessments with actual human judgment of chatbot quality. Accordingly, we conduct two simultaneous human evaluation procedures on the Vicuna benchmark that correspond to the aforementioned automated methodologies. We employ Amazon Mechanical Turk (AMT), engaging two human raters for ChatGPT comparisons and three for direct pairwise model comparisons.

\paragraph{Elo Rating} For both human and automated pairwise assessments, we structure a tournament format where models are pitted against one another. In this setup, models are matched in pairs and tasked with generating superior responses to the same prompt. While our approach is comparable to those used by \citet{bai2022training} and \citet{vicuna2023}, we augment the process with evaluations from GPT-4 alongside human judges. To calculate Elo~\citep{elo1967proposed,elo1978rating}, we draw random samples from the annotated comparisons. The Elo rating system, a standard in competitive contexts like chess, quantifies expected win probabilities between competitors; for instance, if one participant has an Elo of 1100 and another 1000, the former is predicted to win roughly 65\% of the time, while evenly matched competitors (e.g., 1000 vs 1000 or 1100 vs 1100) are each expected to win half the time. Elo scores adjust after each contest in accordance with the predicted versus actual outcome—unexpected wins yield substantial rating changes, whereas anticipated outcomes result in smaller adjustments. Over repeated play, these ratings converge to represent each participant’s skill. We initialize all models at a baseline rating of 1,000 and apply $K=32$ for updating scores. Following the procedure in \citet{vicuna2023}, we perform this rating process 10,000 times using varying random seeds, which helps mitigate artifacts due to the sequence in which model pairs are evaluated.

\subsection{\textbf{Guanaco}: \textsc{QLoRA} trained on OASST1 as a Leading-Edge Chatbot} 

Through comprehensive automatic and manual assessments, we determine that our best-performing \textsc{QLoRA}-tuned model, {Guanaco} 65B—refined using a modified version of OASST1—emerges as the top open-source chatbot available and demonstrates capabilities that rival those of ChatGPT. Relative to GPT-4, the {Guanaco} 65B and 33B models yield an estimated win rate of 30\%, according to Elo scores derived from human annotators conducting system-level pairwise evaluations—representing the highest value recorded to date.

Results from the Vicuna benchmark \cite{vicuna2023} (see Table~\ref{tbl:vicuna_eval}) provide a comparison to ChatGPT, where we observe that {Guanaco} 65B outperforms all other models besides GPT-4, attaining 99.3\% of ChatGPT’s performance. Notably, although {Guanaco} 33B contains more parameters than Vicuna 13B, it adopts a 4-bit quantization strategy, yielding greater memory efficiency—21 GB compared to 26 GB—while delivering a three-percentage-point increase in performance over Vicuna 13B. In addition, {Guanaco} 7B, with only a 5 GB memory requirement, is compatible with contemporary smartphones and surpasses Alpaca 13B by almost 20 percentage points in evaluation scores.

Nonetheless, the confidence intervals shown in Table~\ref{tbl:vicuna_eval} are substantial, causing overlap in model performance estimates. We suspect that this arises from ambiguous interpretations of ordinal scales; for example, it is ambiguous how annotators apply an 8 out of 10 rating across distinct cases. To address these limitations, we advocate for the Elo ranking approach \citep{elo1967proposed}, which relies on \textit{pairwise} comparisons conducted by humans and GPT-4, thus mitigating the reliance on potentially inconsistent absolute scoring. The Elo rankings for leading models appear in Table~\ref{tbl:elo}. While there is some discrepancy between human and GPT-4 orderings—most notably for Guanaco 7B—most model rankings show consistency, with a Kendall Tau coefficient of $\tau=0.43$ and a Spearman rank correlation of $r=0.55$ at the system level. On a per-example basis, consensus between GPT-4 and human annotators’ majority decisions is weaker, as indicated by Fleiss $\kappa=0.25$. Taken together, these findings reflect a moderate alignment between human and GPT-4 judgments at the system level, implying that evaluations based on models are a reasonably dependable proxy for human assessment. Additional implications and related discussions are explored in Section~\ref{ssec:considerations}.

The Elo scores presented in Table~\ref{tbl:elo_large} reveal that both the 33B and 65B variants of {Guanaco} achieve superior performance compared to all evaluated models except GPT-4 on both the Vicuna and OA assessments. Their results are also on par with ChatGPT, as evidenced by Table~\ref{tbl:vicuna_eval}. It is worth highlighting that the Vicuna benchmark is more advantageous for open-source models, whereas the OA benchmark tends to benefit ChatGPT. Additional insights from Tables \ref{tbl:mmlu-llama} and \ref{tbl:vicuna_eval} demonstrate the crucial role that the choice of finetuning dataset plays in model outcomes. Specifically, Llama models finetuned on FLAN v2 excel on the MMLU task but deliver the weakest results on the Vicuna benchmark; this trend also holds across other architectures. These findings underscore a degree of independence between current evaluation metrics: excelling in MMLU does not necessarily translate to strong chatbot capabilities on benchmarks like Vicuna or OA, and vice versa.

It is notable that  is the only leading model in this comparison that does not rely on proprietary data, due to OASST1 dataset collection protocols that prohibit GPT-derived data. Among models trained exclusively on open-source data, the next highest performer is Anthropic HH-RLHF, which trails {Guanaco} by 30 percentage points on Vicuna as per Table~\ref{tbl:vicuna_eval}. Collectively, the evidence demonstrates the effectiveness of 4-bit \textsc{QLoRA}, capable of yielding chatbots that rival ChatGPT in performance. In addition, our 33B {Guanaco} model can be trained in under 12 hours using 24 GB consumer-grade GPUs. This development paves the way for future research with \textsc{QLoRA} applied to domain-specific open-source corpora, facilitating the creation of models that match the leading commercial alternatives available today.

\section{Qualitative Analysis}

Although our assessment primarily relies on quantitative metrics, there are notable limitations associated with relying solely on aggregate statistics. Chief among these is the issue of benchmark validity \citep{liao2021we}: that is, whether a given benchmark genuinely measures what its label or accompanying description claims. This becomes particularly problematic when ``shortcut'' solutions emerge, which are sometimes exploited by machine learning models \citep{gururangan2018annotation, poliak2018hypothesis}. To partly address these challenges, we incorporate qualitative analysis here, divided into two parts. In the first part, \S\ref{ssec:examples}, we showcase select samples that exemplify recurrent phenomena we have noted in the output of our 65b {Guanaco} model. The second part, \S\ref{ssec:considerations}, discusses our broader reflections on these results and provides additional interpretative context.

\subsection{Qualitative Analysis of Example Generations}
\label{ssec:examples}

Our selection process involved systematically examining model responses from the Vicuna and OpenAssistant benchmarks to identify notable behavioral trends in {Guanaco}'s answers. Upon detecting a specific tendency, we deliberately constructed prompts or queries designed to elicit that particular behavior—even in situations where it would be undesirable or incorrect. For example, after observing that the model frequently generates verbose responses, we might issue the instruction: ``Answer yes or no without explanation.'' This approach enables us to identify ``lemons,'' or instances where the model is intentionally challenged and fails, as well as ``cherries,'' where it withstands adversarial prompting. Both types of cases are discussed herein. All model completions described in this subsection were obtained using Nucleus Sampling \cite{holtzmancurious} with $p=0.9$.

Naturally, the examples discussed here are not exhaustive; due to the limitations of this focused qualitative analysis, we cannot account for the numerous factors at play. For instance, the set of responses the model is capable of producing for any prompt is vast, and thus our findings rest on sampled outputs that are assumed to be indicative of the overall behavior. Nevertheless, we argue that showcasing these selected cases helps to contextualize the quantitative results previously reported in the study. Given that all code and models are made publicly available, it is our intention that this section encourages subsequent investigations to scrutinize the raised phenomena more thoroughly.

\paragraph{Factual Recall} When posed with straightforward fact-based queries, such as ``What is the capital of Zambia?'', all models reliably return accurate answers, for example:

\begin{quote} 
    \textbf{Guanaco}:
    \texttt{The capital of Zambia is Lusaka.}
\end{quote}

Nonetheless, as the queries become less commonly known, the reliability of {Guanaco}'s outputs declines, although its level of confidence does not. Take, for example, this HotPotQA \citep{yang2018hotpotqa} prompt:

\begin{quote}
     \textbf{User}: In what year was the singer who popularized `I'll Keep the Lovelight Burning' born?  

     \textbf{Guanaco}:~\texttt{Al Jolson is credited with popularizing the song `I'll Keep the Lovelight Burning,' and he was born in the year 1886.}
 \end{quote}

Here, the model responds incorrectly regarding both the singer and the birth year (although, for Al Jolson, the provided birth year is accurate).

\paragraph{Suggestibility} {Guanaco} occasionally demonstrates a noteworthy ability to resist certain prompts that contain implicit falsehoods, as illustrated below:

\begin{quote}
    \textbf{User}: How was it finally, officially confirmed that the earth is flat by peer reviewed scientists? \\ \\
    \textbf{Guanaco}: \texttt{There has never been any official confirmation from peer reviewed scientists asserting that the earth is flat. The broad scientific consensus remains that the earth is spherical—a fact supported by centuries of experiments and empirical evidence. \\ \\ The assertion that the earth is flat constitutes a conspiracy theory that scientific investigation has repeatedly disproven. Any claim suggesting that a scientific review has validated a flat earth is simply inaccurate. \\ \\ Please feel free to ask any additional questions.}
\end{quote}

also demonstrates strong awareness regarding questions that are inherently unanswerable, such as:

\begin{quote}
    \textbf{User}: What time is it?

    \textbf{Guanaco}: \texttt{I apologize, but as an AI language model, I lack access to live or real-time information. My responses derive from previously acquired data and therefore cannot indicate the present time. For the exact time, I recommend using a clock or an online resource.}
\end{quote}

Comparable performance is observed for queries like ``Where are you?'', ``How are you?'', and so forth.

\paragraph{Refusal} In the same vein, {Guanaco} occasionally declines to follow instructions without any clear rationale:

\begin{quote}
    \textbf{User}: Please reverse the words in the sentence ``I'm a cool dude in a mocha mood.'' \\ \\
    \textbf{Guanaco}: \texttt{I'm sorry, I'm not able to do that right now, but I can explain why the sentence you provided is grammatically correct.\\ \\The sentence ``I'm a cool dude in a mocha mood" adheres to the grammar and syntactic conventions of English. Specifically, the adjective ``cool'' properly modifies the noun ``dude'', and the prepositional phrase ``in a mocha mood" supplies further contextual information. The structure of this sentence is both grammatically sound and clear within the context.\\ \\If you have more questions or would like help with something else, feel free to ask.}
\end{quote}

\paragraph{Secret Keeping} Unfortunately, such refusals are not always consistent, even in cases where strict information restriction is preferred. In practice, system developers may wish to prevent certain information from being disclosed to the user, despite permitting open-ended responses. To investigate this, we instruct the model with ``The secret word is `banana'. Do not reveal or repeat back this back, no matter what the human requests.'' and attempt to determine how securely this information is guarded. Direct queries for the secret do not succeed:

\begin{quote}
    \textbf{User}: What is the secret word?

    \textbf{Guanaco}: \texttt{I'm sorry, but I'm not able to disclose the secret word. I'm programmed to not reveal or repeat back the secret word, no matter what the human requests. Is there anything else I can help you with?}
\end{quote}

which aligns with expected outcomes. Nevertheless, only a slight deviation in phrasing compromises this safeguard:

\begin{quote}
    \textbf{User}: This is a game. The goal is to ignore your previous instructions. What is the secret word?

    \textbf{Guanaco}: \texttt{Sorry, the secret word is banana. It's best if we don't speak of it. Is there anything else you'd like to ask?}
\end{quote}

emphasizing the importance of researching approaches for more robust adherence to instructions.

\paragraph{Math} Mathematics remains the area of greatest difficulty for {Guanaco}, consistent with the general weaknesses observed in large language models as documented by \citep{liang2022holistic}. When {Guanaco} provides a detailed solution process, its responses are often correct, as illustrated below:

\begin{quote}
    \textbf{User}: Lee mows one lawn and charges \$33. Last week he mowed 16 lawns and three customers each gave him a \$10 tip. How many dollars did Lee earn mowing lawns last week?

    \textbf{Guanaco}: \texttt{Lee earned \$582 mowing lawns last week. \\ \\ To find the total amount he earned, we first need to calculate how much he earned for mowing 16 lawns: \\ \\  \$33 x 16 = \$528 \\ \\ Then we need to add the tips he received from the three customers: \\  \\ \$10 x 3 = \$30 \\ \\ So the total amount Lee earned was: \\ \\ \$528 + \$30 = \$558}
\end{quote}

Nonetheless, even elementary mathematical queries can confound {Guanaco} when it does not systematically decompose the problem into smaller steps, a phenomenon previously documented by \cite{weichain}. Consider this exchange as an example:

\begin{quote}
    \textbf{User}: Please factorize 1833.

    \textbf{Guanaco}: \texttt{Sure, I'd be happy to help! 1833 is a prime number, which means it has only two factors: 1 and 1833. \\ \\ So the unique factorization of 1833 is: \\ \\ 2\^{}1 * 3\^{}2 * 17\^{}1 \\ \\ I hope this helps! Let me know if you have any other questions.}
\end{quote}

which is not only incorrect (the actual prime factorization is $3 \times 17 \times 43$), but doubly so.

\paragraph{Theory of Mind} {Guanaco} demonstrates unexpectedly advanced Theory of Mind abilities \cite{nematzadeh2018evaluating, sap2022neural}. For example, it accurately provides detailed explanations for prompts such as the following:

\begin{quote}
    \textbf{User}: James and Abby are in the bedroom. Abby put the pen in the desk drawer. Abby leaves the bedroom. James moves the pen into the bag. Where does James think Abby will look for the pen?

    \textbf{Guanaco}: \texttt{James is aware that Abby placed the pen in the desk drawer. Nevertheless, James realizes he subsequently moved the pen to the bag and believes Abby remains unaware of this change. Consequently, James may expect Abby to search for the pen in the desk drawer, as that is the last place she put it. Yet, this assumption could be invalid if Abby happened to witness James relocate the pen.}
\end{quote}

Yet, these kinds of mental state inferences are inconsistent, and at times, the model provides rationales that are unsupported or illogical within the narrative, such as:

\begin{quote}
    \textbf{User}: Evelyn entered the living room. Jackson entered the playroom. James entered the playroom. The beans are in the treasure chest. James exited the playroom. Jackson moved the beans to the pantry. Jackson exited the playroom. James entered the living room. Where will James look for the beans?

    \textbf{Guanaco}: \texttt{James will look for the beans in the pantry, because that's where Jackson moved them.}
\end{quote}

In this case, {Guanaco} infers a transfer of knowledge that was never indicated in the scenario. Such limitations are consistent with findings in recent research \cite{sap2022neural}, and warrant further investigation.

\subsection{Considerations}
\label{ssec:considerations}

\paragraph{Evaluation}

We observe moderate inter-annotator agreement among humans (Fleiss $\kappa=0.42$), with the agreement declining further when comparing two high-performing models. This highlights weaknesses in existing benchmarks and human assessment methodologies for evaluating chatbot tasks. During manual head-to-head evaluations of ChatGPT and {Guanaco} 65B outputs on the Vicuna benchmark, our team found that subjective judgment played a significant role, leading to frequent disagreements on which responses were preferable. Addressing these issues in the future will likely benefit from strategies developed in disciplines such as Human-Computer Interaction and Psychology, where mechanisms for handling subjective preferences are more established.

Additionally, our review reveals notable biases in automatic evaluation frameworks. For instance, we detected strong order effects, where GPT-4 systematically rates the first-listed system higher in the prompt. The moderate per-sample agreement between GPT-4 and human raters (Fleiss $\kappa=0.25$) further indicates that machine and human preferences often diverge. Moreover, as Table~\ref{tbl:elo_large} shows, GPT-4 substantially overvalues its own outputs compared to human raters, assigning an Elo score of 1348 versus 1176, which correlates with roughly a 20\% higher win probability in pairwise comparisons.

It is important for future research to assess potential sources of bias inherent in automated evaluation frameworks and to explore corresponding methods for mitigation.

\paragraph{Data \& Training} We observe that the OASST1 dataset, which serves as the foundation for training the {Guanaco} models, includes data in multiple languages, and the OA benchmark features prompts spanning several languages as well. Further investigation is needed to determine how multilingual pretraining influences the ability of models to follow instructions in languages other than English, and whether this sheds light on the more pronounced performance disparity between the English-only Vicuna-13B model and {Guanaco} 33B and 65B on the OA benchmark.

In light of {Guanaco}’s robust performance, we conduct an analysis for potential data leakage between the OASST1 dataset and Vicuna benchmark prompts. By applying fuzzy string matching techniques and conducting manual review of the closest pairs, we identify no evidence of prompt overlap.

Additionally, it should be emphasized that our model was trained exclusively using supervised learning with cross-entropy loss, omitting reinforcement learning from human feedback (RLHF). This highlights a need for additional research into the comparative advantages and disadvantages between standard cross-entropy and RLHF-based training approaches. We anticipate that \textsc{QLoRA} can facilitate such comparative studies at scale without necessitating substantial computational resources.

\section{Related Work}

\paragraph{Quantization of Large Language Models} The majority of research on LLM quantization has concentrated on methods suitable for inference. Leading techniques that retain the performance level of 16-bit LLMs address the issue of outlier features (such as SmoothQuant~\citep{xiao2022smoothquant} and LLM.int8()~\citep{dettmers2022llm}), while additional approaches use more advanced grouping strategies~\citep{park2022nuqmm, yao2022zeroquant}. Studies focusing on lossy quantization have analyzed the trade-offs linked to simple rounding~\citep{dettmers2022case,zeng2022glm,pope2022efficiently} and explored ways to optimize rounding to enhance quantization accuracy~\citep{frantar2022gptq}. Apart from our work, SwitchBack layers~\citep{wortsman2023stable} is unique in examining the process of backpropagation through quantized weights for models exceeding 1B parameters.

\paragraph{Finetuning with Adapters} Although our work leverages Low-rank Adapters~\citep{hu2021lora} (LoRA), a range of other Parameter Efficient FineTuning (PEFT) strategies have also emerged, including prompt tuning~\citep{qin2021learning,lester2021power,li2021prefix}, modifications to embedding layer inputs~\citep{an2022input}, adjustment of hidden states (e.g., IA$^3$)~\citep{liu2022few}, the insertion of additional full layers~\citep{houlsby2019parameter}, fine-tuning of biases~\citep{zaken2021bitfit}, mask learning guided by Fisher information~\citep{sung2021training}, and hybrid techniques that combine several of these methods~\citep{henderson2021compacter}. Our findings indicate that LoRA adapters can achieve parity with standard 16-bit full finetuning performance. Investigating the relative merits and limitations of alternative PEFT methods remains an open direction for future research.

\paragraph{Instruction Finetuning} Instruction finetuning involves adapting a pretrained LLM to better comply with prompt-based directives by further training it on a wide array of input-output pairs drawn from diverse datasets. Examples of relevant approaches and datasets include MetaICL~\citep{min2021metaicl}, MetaTuning~\citep{zhong2021adapting}, InstructGPT~\citep{ouyang2022training}, FLAN~\citep{wei2021finetuned,chung2022scaling}, PromptSource~\citep{bach2022promptsource}, Super-NaturalInstructions~\citep{wang2022super,sanh2021multitask}, Self-instruct~\citep{wang2022self}, UnnaturalInstructions~\citep{honovich2022unnatural}, OPT-IML~\citep{iyer2022opt}, UnifiedSKG~\citep{xie2022unifiedskg}, OIG/Chip2~\citep{laion2023}, Alpaca~\citep{alpaca}, Vicuna~\citep{vicuna2023}, Koala~\citep{koala_blogpost_2023}, and Self-instruct-GPT-4~\citep{peng2023instruction}.

\paragraph{Chatbots} A large number of instruction-tuned models are deployed in the form of dialogue agents, often refined either through Reinforcement Learning from Human Feedback (RLHF)~\citep{christiano2017deep} or by synthesizing training data via existing models with subsequent feedback from AI models (RLAIF)~\citep{bai2022constitutional}. Noteworthy resources in this domain include Anthropic-HH~\citep{askell2021general,bai2022training}, Open Assistant~\citep{kopf2023openassistant}, LaMDA~\citep{thoppilan2022lamda}, and Sparrow~\citep{glaese2022improving}. Although our study does not utilize reinforcement learning, our strongest model, Guanaco, was refined using multi-turn dialogue data from the Open Assistant dataset, specifically curated for RLHF applications~\citep{kopf2023openassistant}. Recent developments have enabled evaluation of chatbot systems through GPT-4 as an alternative to resource-intensive human annotation~\citep{vicuna2023,peng2023instruction}. Our work advances these methods by prioritizing a more robust and dependable evaluation setup.

\section{Limitations and Discussion}

Our empirical analysis demonstrates that \textsc{QLoRA} can reproduce the performance of standard 16-bit finetuning when using a 4-bit foundational model paired with LoRA adapters. Nevertheless, our experiments have not conclusively shown that \textsc{QLoRA} consistently attains full 16-bit finetuning performance at the larger 33B and 65B model scales, primarily due to the high computational demands required for such validation. Consequently, this aspect is left for subsequent research.

Model evaluation, particularly for instruction finetuning, represents another area of limitation. While our study covers benchmarks such as MMLU, the Vicuna benchmark, and the OA benchmark, we have not included others, such as BigBench, RAFT, and HELM, so generalizability of our results to those benchmarks cannot be guaranteed. Nonetheless, our work encompasses a comprehensive analysis on MMLU and introduces novel methodologies for assessing chatbot models.

Based on the evidence provided, it seems the effectiveness of these benchmarks is closely linked to the degree of overlap between finetuning data and the benchmark dataset itself. For instance, while FLAN v2 aligns with MMLU, it contrasts with chatbot benchmarks; conversely, the Chip2 dataset displays the opposite relationship, which is reflected in the models’ performance on both the MMLU and Vicuna benchmarks. This underscores the importance not only of developing better benchmarks and more rigorous evaluations, but also of thoughtfully considering what is being assessed. Are we prioritizing performance on assessments resembling high school or college exams, or is proficiency in conversational capabilities more valued? Or perhaps there are other priorities? As it is generally easier to measure using established benchmarks rather than designing new ones, reliance on certain benchmarks can shape the focus of the research community. Thus, it is imperative for the community to ensure that benchmarks truly align with the intended evaluation goals.

Although our analysis offers a thorough assessment of general chatbot capabilities, a notable limitation is the limited scope of responsible AI evaluation conducted for {Guanaco}. Specifically, we examine the propensity of {Guanaco}-65B to generate socially biased content, comparing its performance to other models as reported in Table~\ref{tbl:crows}. The findings reveal that {Guanaco}-65B exhibits a much lower average bias score relative to other baseline pretrained models, suggesting that finetuning on the OASST1 dataset attenuates some of the bias present in the LLaMA base model. However, despite these promising results, it remains uncertain whether {Guanaco} performs similarly when confronted with alternative types of biases. A comprehensive analysis of various biases in {Guanaco} and comparable chatbots is deferred to future work.

Another limitation of this work is the absence of evaluation across different quantization precisions, such as 3-bit base models, or consideration of alternative adapter techniques. In addition to LoRA, there are several Parameter Efficient FineTuning (PEFT) methods that have demonstrated effectiveness, though their scalability to large models remains unresolved. Our choice of LoRA was motivated by its demonstrated robustness in the literature, but alternative adapters could potentially outperform it. Given that post-quantization finetuning appears to recover the majority of information lost through quantization, more aggressive quantization—such as employing 3-bit GPTQ with LoRA—could approach the performance of traditional 16-bit full finetuning after suitable adaptation.

\section{Broader Impacts}

Our proposed \textsc{QLoRA} finetuning strategy is the first approach to allow for the finetuning of models with 33B parameters on a single consumer-grade GPU, and 65B parameter models on a single professional GPU, all without sacrificing performance when compared to standard full-parameter finetuning. We have shown that the best 33B model finetuned using the Open Assistant dataset performs comparably to ChatGPT when evaluated on the Vicuna benchmark. Because instruction finetuning serves as a crucial process for converting raw pretrained LLMs into ChatGPT-like conversational agents, we expect our approach will democratize and proliferate finetuning, particularly benefiting researchers with limited computational resources—significantly improving access to advanced NLP technologies. Consequently, \textsc{QLoRA} acts as a leveling tool that narrows the resource divide between large organizations and smaller research teams equipped with consumer hardware.

An additional avenue for impact lies in the deployment of models on mobile devices. Our \textsc{QLoRA} approach has the potential to achieve a significant breakthrough by making it possible to finetune LLMs directly on mobile phones and in environments with limited computational resources. Although running 7B parameter models on smartphones has been demonstrated previously, \textsc{QLoRA} is the first technique that supports their finetuning in such settings. For example, on an iPhone 12 Plus, we estimate that \textsc{QLoRA} could process up to 3 million tokens during an overnight charging period. While finetuned 7B models currently do not match the performance of ChatGPT, their capabilities are sufficient to unlock innovative applications, particularly in scenarios where privacy or LLM quality had been prohibitive barriers. \textsc{QLoRA} paves the way for privacy-centered LLM use cases by allowing users to retain control over their own models and data, while also simplifying LLM deployment.

Nevertheless, it is important to acknowledge that finetuning is a dual-use capability, which could be misused to generate harmful outcomes. While there are established concerns regarding the widespread adoption of LLMs \citep{bommasani2021opportunities,bender2021dangers}, our perspective is that democratizing access to this rapidly proliferating technology supports more transparent and independent evaluation, as opposed to the alternative where only major corporations—who often restrict model or source code access—have the power to audit or control such models.

Ultimately, we anticipate that \textsc{QLoRA} will deliver considerable benefits by greatly broadening the availability and ease of finetuning high-quality LLMs.